\section{Introduction}
\label{sec:introduction}

Manufacturing failures, constrained inputs, discontinuations, demand shocks, and slow recovery can all contribute to drug shortages. FDA analyses emphasize quality problems, limited capacity, and small supplier bases \cite{fda2019shortages}. A public shortage record may identify a presentation, company, availability state, and reported reason, but analysts must consult separate regulatory sources to interpret quality risk, supplier control, or change management.

Those sources are difficult retrieval targets. A short clause may inherit scope from several parent headings, and PDF conversion can separate a table value from its conditions. Guideline numbers, defined phrases, and product identifiers often carry more evidential weight than broad semantic similarity. Structured facts create another requirement: NDCs, manufacturers, active ingredients, and events form relations, not citable prose. Treating a graph edge as a quotation or a similar passage as a database fact obscures provenance.

We address this problem as evidence retrieval rather than answer generation. The framework returns source passages with hierarchy and separately returns provenance-bearing graph paths. It compares lexical, dense, enriched, hierarchical, fused, and neural rankings under frozen evaluation protocols. The final evidence set was independently labelled by the author and a domain reviewer, jointly adjudicated, and executed once after method locking.

The contribution is empirical as well as architectural. We test where hierarchy and graphs help, where they do not, and whether a biomedical retrieval model transfers to pharmaceutical regulation. The system does not predict shortage duration or patient harm. It exposes what a frozen source says and how a structured record connects named entities.
